% Generated draft: requires mathematical and source review.
\documentclass[11pt]{article}
\usepackage[margin=25mm]{geometry}
\usepackage{fontspec}
\setmainfont{Latin Modern Roman}
\usepackage{amsmath,amssymb,mathtools,unicode-math}
\setmathfont{Latin Modern Math}
\usepackage{enumitem,needspace,xurl,hyperref}
\setlength{\parindent}{0pt}
\setlength{\parskip}{5pt}
\setlength{\emergencystretch}{3em}
\newcommand{\N}{\mathbb{N}}
\newcommand{\Z}{\mathbb{Z}}
\newcommand{\Q}{\mathbb{Q}}
\newcommand{\R}{\mathbb{R}}
\newcommand{\C}{\mathbb{C}}
\newcommand{\F}{\mathbb{F}}
\newcommand{\A}{\mathbb{A}}
\newcommand{\PP}{\mathbb{P}}
\newcommand{\E}{\mathbb{E}}
\newcommand{\eps}{\varepsilon}
\newcommand{\dd}{\,\mathrm{d}}
\begin{document}
\section*{UnsolvedMath ID: 700912}
\textbf{Source problem code:} catalogue entry
\subsection{Definitions and mathematical statement}
Let \(p\) be a prime. A pro-\(p\) group is a topological group isomorphic to an inverse limit of finite \(p\)-groups. It is finitely generated if some finite subset generates a dense subgroup. For each integer \(r\geq 1\), let \(F_r=F(X_1,\ldots,X_r)\) be the free group on \(r\) generators. A non-trivial word means an element of \(F_r\) different from its identity element; evaluating a word in a group means substituting group elements for its generators.

The proposed assertion is the following. For every prime \(p\) and every finitely generated pro-\(p\) group \(G\), suppose there exist an integer \(n\geq 1\), elements \(g_1,\ldots,g_n\in G\), an open subgroup \(H\leq G\), and a word \(w\in F_n\setminus\{1\}\) such that
\[
  w(g_1h_1,\ldots,g_nh_n)=1_G
  \qquad\text{for every }(h_1,\ldots,h_n)\in H^n.
\]
This condition is called a coset identity. The subgroup \(H\) is not assumed to be normal. Must there then exist an integer \(m\geq 1\) and a word \(v\in F_m\setminus\{1\}\) such that
\[
  v(b_1,\ldots,b_m)=1_G
  \qquad\text{for every }(b_1,\ldots,b_m)\in G^m?
\]
The word \(v\) and its number of variables may depend on \(G\); no uniform identity is requested.

The supplied background states that a positive answer would imply an analogue of the Tits Alternative for finitely generated pro-\(p\) groups, but does not formulate that analogue. It also attributes to J. S. Wilson and E. I. Zelmanov the result that, under the coset-identity hypothesis, the graded Lie algebra \(L_p(G)\) associated with the dimension subgroups of \(G\) in characteristic \(p\) satisfies a polynomial identity. These are supplied contextual claims, not conclusions established in this draft.

A supplied comment dated 2017 reports a positive answer for finitely generated linear groups, citing M. Larsen and A. Shalev. The comment does not specify the field, characteristic, or the intended interpretation of finite generation in that linear-group formulation; no more precise formulation of that special case is asserted here.
\subsection{Short English statement}
If a finitely generated pro-\(p\) group satisfies a non-trivial group word on every tuple from specified cosets of one open subgroup, must it satisfy some non-trivial group word on all tuples from the whole group? The supplied registry labels this problem open and includes a 2017 comment reporting a positive result for finitely generated linear groups.
\subsection{Sources}
\begin{itemize}
\item Local UnsolvedMath registry, supplied complete record: permanent ID \(700912\), exact source code \texttt{catalogue entry}, title \emph{Kourovka Notebook Problem 12.95}. The supplied status is \texttt{open}; this label has not been independently verified.
\item A. Shalev, proposer, 1992. \emph{The Kourovka Notebook (Unsolved Problems in Group Theory)}, original problem 12.95, 12th issue (1992). The supplied background identifies the consulted edition as the 21st issue edition (2026), \texttt{arXiv:1401.0300v48}. The registry describes the collection as edited by E. I. Khukhro and V. D. Mazurov.
\item J. S. Wilson and E. I. Zelmanov, \emph{J. Pure Appl. Algebra} \textbf{81} (1992), 103--109. Supplied attribution for the polynomial-identity result concerning \(L_p(G)\); no article title or link is supplied.
\item M. Larsen and A. Shalev, \emph{Algebra Number Theory} \textbf{10}, no. 6 (2016), 1359--1371. Cited by the supplied \emph{Comment of 2017} for the positive result for finitely generated linear groups; no article title or link is supplied.
\end{itemize}
\end{document}
